\paragraph{\trKodaira--\trSpencer{}映射。}\label{sub:kodaira-spencer}
设 $S$ 是一个带\trStriplikeEnd{}的\trRiemann\trSurface。在\trTangentBundle{}上取标准的 $\bar{\partial}$-算子，边界
条件取在\trTotallyReal\trSubbundle{}$T(\partial S)$ 之中。用\cref{sub:cauchy-riemann-operators} 的话说，这是\trAdmissible{}的（但并非\trNondegenerate），
从而在取小\trWeight{}$\epsilon>0$ 的\trWeighted\trSobolevSpace{}上成为\trFredholm{}的，见
\eqref{eq:cauchy-riemann-fredholm}。实际上，我们想容许一个稍大一些的\trDomain
\begin{equation}\label{eq:enlarged-domain}
W^{k,p;\epsilon}(S,TS,T(\partial S))^{+}\supset W^{k,p;\epsilon}(S,TS,T(\partial S)),
\end{equation}
它由那些在每条\trStriplikeEnd{}上形如 $\epsilon_\zeta^*Z=c_\zeta\partial_s+\widetilde{Z}_\zeta$（其中 $c_\zeta\in\RR$、
$\widetilde{Z}_\zeta\in W^{k,p;\epsilon}$）的\trVectorField{}组成。这增添了有限多个维数（每条端增添一个，
对应于诸常数 $c_\zeta$）；因此，$\bar{\partial}$ 仍扩张为一个\trFredholmOperator
\begin{equation}\label{eq:dbar-plus-fredholm}
\bar{\partial}^{+}:W^{k,p;\epsilon}(S,TS,T(\partial S))^{+}
\longrightarrow W^{k-1,p;\epsilon}(S,\Omega^{0,1}_{S}\otimes TS).
\end{equation}
把这个算子的\trCokernel{}记作 $\mathscr{T}_{S}$；它既与 $(k,p)$ 无关，也与（适当小的）$\epsilon>0$ 无关。
若 $S$ 是稳定的（即没有\trAutomorphism{}的连续\trFamily），则\eqref{eq:dbar-plus-fredholm} 是%
\trInjective{}的，且 $\mathscr{T}_{S}$ 的维数为 $|\Sigma|-3\chi(\hat{S})$。
设 $\scr{S}\to\scr{R}$ 是带\trStriplikeEnd{}的\trRiemann\trSurface\trFamily。取一点 $r_0\in\scr{R}$，记 $S=\scr{S}_{r_0}$ 为其%
\trFiber，并取一个使\trStriplikeEnd{}恒定的局部\trTrivialization{}$\Psi$，如\eqref{eq:family-trivialization} 中
那样。对 $S$ 上由\trPullback{}\trComplexStructure{}构成的\trFamily{}$(\Psi^*I_{\scr{S}})$ 求\trDifferential，便对 $(T\scr{R})_{r_0}$
中的每个切向量给出 $C^{\infty}_{c}(S,\mathrm{End}^{0,1}(TS))=C^{\infty}_{c}(S,\Omega^{0,1}_{S}\otimes TS)$ 中的
一个元素；把它\trProjection{}到 $\mathrm{coker}(\bar{\partial}^{+})$ 之上，就定义了\textbf{\trKodaira--\trSpencer{}映射}%
（Kodaira–Spencer map）
\begin{equation}\label{eq:kodaira-spencer-map}
\kappa_{r_0}:(T\scr{R})_{r_0}\longrightarrow\mathscr{T}_{S}.
\end{equation}
这与\trTrivialization{}的选取无关，而且（尽管这一点从我们的定义里不易立刻看出——因为我们偏爱与%
\trStriplikeEnd{}结构紧密相连的函数空间）也与\trStriplikeEnd{}的选取无关。称 $\scr{S}\to\scr{R}$ 是\textbf{无穷小%
\trMiniversal}（infinitesimally miniversal）的，若\eqref{eq:kodaira-spencer-map} 对每个 $r_0\in\scr{R}$ 都是%
\trBijective{}的。要为这个名称辩护，就得证明这类\trFamily{}具有局部\trVersality。这并不困难，
但我们实际上用不上它。感兴趣的读者可参看带\trMarkedPoint{}的\trRiemann\trSurface{}的代数几何%
\trDeformation{}理论 [72, p. 94]，二者完全类似。
